\paragraph{The rational affine comparison also fails.}
The following calculation settles the rational assertion of literal Lead E10; it concerns $H^1(\mathcal C,\OO)\otimes_WK$, with precisely the old objects and coefficients.

\begin{lemma}[Finite weak Nakayama lemma]\label{lem:DE-weak-finite}
Let $A\to B$ be a homomorphism of weakly complete, finitely generated $W$-algebras.  If $B/pB$ is finite over $A/pA$, then $B$ is finite over $A$.
\end{lemma}
\begin{proof}
Here is a proof, to avoid attributing this finite-module statement to a source that only states the surjectivity case.  Choose finitely many weak generators $y_1,\ldots,y_r$ of $B$ over $A$.  Integrality modulo $p$ supplies a monic polynomial $P_j(T)\in A[T]$ for each $j$, of degree $d_j$, with $P_j(y_j)\in pB$.  Lift the quotient $P_j(y_j)/p$ to a dagger series $G_j$ in the variables $y_1,\ldots,y_r$.  Thus $B$ is a quotient of
\[
 A\langle y_1,\ldots,y_r\rangle^\dagger/
       (P_j(y_j)-pG_j:1\leq j\leq r).
\]
This quotient is generated over $A$ by the finitely many monomials $y^\nu$ with $0\leq\nu_j<d_j$.

For completeness, the convergence needed in that last assertion is as follows.  Present $A$ as a quotient of an integral dagger power-series algebra in finitely many base variables.  Choose overconvergent Banach radii for these variables and for the $y_j$, all strictly greater than $1$ and sufficiently near $1$.  Choose the variable radii sufficiently large relative to the base radius that ordinary successive monic division by the $P_j(y_j)$ is bounded: its remainder operator $R$ has norm at most $1$, and its quotient operators $Q_j$ have norm at most $\rho_j^{-d_j}$.  This is possible because each lower coefficient of $P_j$ is integral and its norm tends to at most $1$ as the base radius tends to $1$.  All $G_j$ converge at the chosen radii.  These radii can be decreased towards $1$ so that
\[
             T(f)=p\sum_jG_jQ_j(f)
             \quad\hbox{has operator norm strictly less than }1.
\]
The division identity is
\[
 f=R(f)+\sum_j(P_j-pG_j)Q_j(f)+T(f).
\]
Iteration gives a convergent Neumann series $\sum_{n\geq0}R(T^nf)$ for the remainder.  Its coefficients converge in the fixed base Banach algebra at radius greater than $1$, hence belong to $A$, and its $y$-degrees remain bounded by the $d_j$.  The corresponding quotient series also converge at these radii.  Thus this is an identity of dagger series, not merely an identity in a $p$-adic completion.  It proves finite generation and hence the lemma.
\end{proof}

\begin{theorem}\label{thm:DE-old-Gm}
For the literal flat site of 7:18,
\[
 H^1(\mathcal C_{\mathbb G_{m,k}},\OO)=0,
 \qquad
 H^1_{\rm MW}(\mathbb G_{m,k}/W)\otimes_WK
       =K\,[dx/x].
\]
Consequently test (a) fails even after tensoring with $K$, and the rational rigid comparison asserted in literal Lead E10 is false.
\end{theorem}
\begin{proof}
Put $A=W\langle x,x^{-1}\rangle^\dagger$.  Its associated object covers the final sheaf of $\mathcal C$: on every object $B$, a lift of the invertible coordinate $x$ is invertible, since the kernel of $B/pB\to\OO(U)$ is nilpotent and $p$ lies in the Jacobson radical of $B$.  The universal property of weak completion then gives $A\to B$.

An additive $\OO$-torsor is trivial on the object $A$.  Indeed its restriction there is governed by Zariski coverings of the affine weak formal scheme, and \cite[\S2, Theorem~14, p.~10]{Meredith1972} gives $H^1=0$ for its structural sheaf.  Equivalently, one refines the local trivializing coverings of that object by principal weak formal opens and applies that theorem.  After choosing a trivialization, the torsor is therefore represented by a natural additive \v Cech cocycle $u$: on an object $B$ with two maps $a_0,a_1:A\to B$, it has value $u_B(a_0,a_1)\in B$, and
\[
 u_B(a_0,a_2)=u_B(a_0,a_1)+u_B(a_1,a_2).
\]
Changing the trivialization subtracts a coboundary $a_1(f)-a_0(f)$, $f\in A$.

First evaluate on the allowed objects
\[
                P_N=A[t]/(t^N),\qquad
                a_0(x)=x,\quad a_1(x)=x+t.
\]
These give $u(x,t)\in A[[t]]$ with $u(x,0)=0$ and
\[
              u(x,t+s)=u(x,t)+u(x+t,s).
\]
The latter equality follows on the objects with two independent nilpotent variables, and hence in the formal power-series ring.  Differentiating in $s$ at $0$ shows that $\partial_tu(x,t)=a(x+t)$, where $a(x)=\partial_tu(x,0)\in A$.  Write
\[
 a(x)=\sum_{m\in\mathbb Z}a_mx^m,\quad
 c=a_{-1},\quad b_s=a_{s-1}/s\quad(s\ne0).
\]
Over $K$, formal integration gives
\[
 u(x,t)=c\log(1+t/x)+
       \sum_{s\ne0}b_s\bigl((x+t)^s-x^s\bigr).
\]
All coefficients of $u$ are integral.  Its coefficient at $t^nx^{-n}$ is $(-1)^{n-1}c/n$, so $c=0$, by taking $n$ to be arbitrarily large powers of $p$.  For each nonzero integer $s$, the coefficient at $t^nx^{s-n}$ is $b_s\binom{s}{n}$.  Take $n=p^{v_p(s)}$.  The binomial coefficient is a $p$-adic unit, also for negative $s$: in $\mathbb F_p[[z]]$, the first nonconstant coefficient of $(1+z)^s$ occurs at $z^{p^{v_p(s)}}$ and is $s/p^{v_p(s)}$ modulo $p$.  Hence $b_s\in W$.  The coefficients $a_m$ have linear $p$-adic decay in $|m|$, whereas division by $s$ loses at most $\log_p|s|$.  Therefore
\[
                         f(x)=\sum_{s\ne0}b_sx^s\in A.
\]
After subtracting its coboundary, $u$ vanishes on every $P_N$.

It remains to check that these evaluations detect the cocycle on every old object; a formal-diagonal calculation alone would not suffice.  This is local on its underlying open $U$, so take $U$ principal in $\mathbb G_m$, and let $A_U$ be the dagger localization of $A$ at a lift of its denominator.  Both maps $a_0,a_1$ extend to $A_U$.  The preceding vanishing also holds on
$P=A_U[t]/(t^N)$, by localization.  Set $\delta=a_1(x)-a_0(x)\in B$.  Its reduction is nilpotent, so, for every integer $m\geq1$, choose $N$ with $\delta^N\in p^mB$.  Evaluation at $x=a_0(x)$ and $t=\delta$ gives a map $P\to B/p^mB$.  Form the ordinary ring fibre product
\[
                         C=B\times_{B/p^mB}P.
\]
Both $B$ and $P$ are finite over $A_U$: modulo $p$, the first is a finite-type nilpotent thickening of $\OO(U)$ with a chosen section, and the second is finite free; apply Lemma~\ref{lem:DE-weak-finite}.  Since $A_U$ is noetherian, $C\subset B\oplus P$ is a finite $A_U$-module.  It is weakly complete, because evaluations of convergent dagger series preserve the defining equality modulo $p^m$; thus it is a weakly complete algebra of finite type over $W$.  It is $W$-flat as a submodule of $B\oplus P$.

The projection $C\to P$ is surjective with kernel $(p^mB,0)$.  The image of this ideal in $C/pC$ has square zero: the product of two such elements is $p$ times an element of $(p^mB,0)\subset C$, since $2m-1\geq m$.  The kernel of $P/pP\to\OO(U)$ is nilpotent.  Consequently $C/pC\to\OO(U)$ has nilpotent kernel, and $C$ is an allowed old-site object.  The pairs $(a_0,x)$ and $(a_1,x+t)$ define two maps $A\to C$.  Naturality of the cocycle and its vanishing on $P$ imply
\[
                         u_B(a_0,a_1)\in p^mB.
\]
This holds for every $m$, and $B$ is $p$-adically separated.  Hence $u_B(a_0,a_1)=0$.  The torsor is trivial, proving the first assertion.

For the second assertion, termwise integration of an overconvergent Laurent series over $K$ removes every differential term except its residue multiple of $dx/x$; the same linear-decay estimate ensures that the primitive remains overconvergent.  The residue is zero on exact forms.  Thus the displayed one-dimensional Monsky--Washnitzer group is obtained directly.  Its standard identification with rigid cohomology for a smooth affine variety, or the elementary annulus calculation itself, shows that the rational old-site group also differs from rigid cohomology.
\end{proof}

The smooth proper assertion (b) is not settled by this affine counterexample.  It remains a separate gate in the present draft.
